    % !TEX root = ../main.tex
\chapter{Ultrashort-Pulse Intuition}
\label{app:phase_kick_intuition}

This appendix collects an optional geometric intuition for the role of pulse carrier phases in pathway selection. It complements the formal treatment in~\autoref{subsec:spctr_phase_matching_condition,subsec:spctr_rephasing_nonrephasing,sec:spctr_phase_cycling} without changing the main derivation used in the thesis.

\section{Carrier Phase, Pulse Area, and Bloch-Sphere Picture}
\label{app:sec:phase_kick_pulse_area_intuition}

The carrier phases $\phi_j$ in~\autoref{eq:spctr_E_three_pulses_space} are experimentally controlled parameters that imprint well-defined phase factors on driven coherences and hence on the nonlinear polarisation. This provides the pathway selection used throughout this thesis.
This geometric picture is used as intuition. All quantitative results in this thesis are obtained from the full, numerically exact propagation and the phase-selection procedures developed in~\autoref{sec:spctr_phase_cycling}.

To connect this intuition directly to the main text, start from the scalar semiclassical light--matter Hamiltonian in~\autoref{eq:spctr_semiclassical_hamiltonian_scalar} and specialize it to the two-level reference system introduced in~\autoref{sec:spctr_two_level_system_linear_response}. The aim of this appendix section is only to show how the pulse-resolved quantity $\Omega_j$ arises from the same dipole coupling already used in the formal response-function derivation.

\subsection{From the Dipole Hamiltonian to the Complex Rabi Frequency}

For the two-level system, the field-free Hamiltonian and dipole operator are
\begin{equation*}
    \hat{H}_{0} = \omega_{\mathrm{eg}}\,\ket{\mathrm{e}}\!\bra{\mathrm{e}},
    \qquad
    \hat{\mu}
    = \mu_{\mathrm{ge}}\,\ket{\mathrm{g}}\!\bra{\mathrm{e}} + \mu_{\mathrm{eg}}\,\ket{\mathrm{e}}\!\bra{\mathrm{g}},
\end{equation*}
as given in~\autoref{eq:spctr_tls_hamiltonian,eq:spctr_tls_dipole}. Inserting these into~\autoref{eq:spctr_semiclassical_hamiltonian_scalar} gives
\begin{equation*}
    \hat{H}_{\mathrm{S}}(t)
    = \omega_{\mathrm{eg}}\,\ket{\mathrm{e}}\!\bra{\mathrm{e}}
    - E(t)\Big(\mu_{\mathrm{ge}}\,\ket{\mathrm{g}}\!\bra{\mathrm{e}} + \mu_{\mathrm{eg}}\,\ket{\mathrm{e}}\!\bra{\mathrm{g}}\Big).
\end{equation*}
To isolate the effect of one pulse, insert the $j$th contribution from~\autoref{eq:spctr_E_three_pulses_space},
\begin{equation*}
    E_j^{(+)}(t)
    = E_j(t-t_j)\,e^{-\mathrm{i}\omega_{\mathrm{L}} t + \mathrm{i}\phi_j},
    \qquad
    E_j^{(-)}(t)=\bigl[E_j^{(+)}(t)\bigr]^*,
\end{equation*}
so that $E(t)=E_j^{(+)}(t)+E_j^{(-)}(t)$ for that pulse. In the interaction picture with respect to $\hat{H}_{0}$, the transition operators acquire the phases
\begin{equation*}
    e^{+\mathrm{i}\hat{H}_0 t}\ket{\mathrm{e}}\!\bra{\mathrm{g}}e^{-\mathrm{i}\hat{H}_0 t}=e^{+\mathrm{i}\omega_{\mathrm{eg}}t}\ket{\mathrm{e}}\!\bra{\mathrm{g}},
    \qquad
    e^{+\mathrm{i}\hat{H}_0 t}\ket{\mathrm{g}}\!\bra{\mathrm{e}}e^{-\mathrm{i}\hat{H}_0 t}=e^{-\mathrm{i}\omega_{\mathrm{eg}}t}\ket{\mathrm{g}}\!\bra{\mathrm{e}}.
\end{equation*}
The interaction Hamiltonian therefore contains terms oscillating at $\omega_{\mathrm{eg}}\pm\omega_{\mathrm{L}}$. Near resonance, the RWA keeps only the slowly varying combinations $\ket{\mathrm{e}}\!\bra{\mathrm{g}}\,E_j^{(+)}$ and $\ket{\mathrm{g}}\!\bra{\mathrm{e}}\,E_j^{(-)}$, while the counter-rotating terms are discarded. Up to the sign and factor conventions associated with the chosen rotating frame, this yields
\begin{equation*}
    \hat{H}_{\mathrm{int},j}(t)
    \;\propto\;
    \mu_{\mathrm{eg}}E_j(t-t_j)e^{+\mathrm{i}\phi_j}\,\ket{\mathrm{e}}\!\bra{\mathrm{g}}
    + \mu_{\mathrm{ge}}E_j^*(t-t_j)e^{-\mathrm{i}\phi_j}\,\ket{\mathrm{g}}\!\bra{\mathrm{e}}.
\end{equation*}
This motivates defining the complex Rabi frequency of pulse $j$ as the slowly varying dipole--field coupling amplitude,
\begin{equation*}
    \Omega_j(t) \equiv \mu_{\mathrm{eg}}E_j(t).
\end{equation*}
Depending on whether the carrier phase is attached to the field envelope or absorbed into the definition of the rotating-frame basis, one may equivalently write $\phi_j\to -\phi_j$. Using the convention adopted for the Bloch-sphere picture below, the interaction Hamiltonian in the laser-rotating frame can be expressed as
\begin{equation*}
    \hat{H}_{\mathrm{int},j}(t)
    \;\propto\;
    \Omega_j(t-t_j)
\Big( e^{-\mathrm{i}\phi_j}\,\ket{\mathrm{e}}\!\bra{\mathrm{g}} + e^{+\mathrm{i}\phi_j}\,\ket{\mathrm{g}}\!\bra{\mathrm{e}} \Big).
\end{equation*}
This is the form in which the carrier phase appears as a rotation angle in the equatorial plane of the Bloch sphere.
Changing $\phi_j$ rotates the complex drive while leaving the pulse envelope unchanged. On the Bloch sphere this corresponds to a rotation about an equatorial axis with azimuth angle $\phi_j$ (cf.~\cite{grolletal2025fundamentalsheterodynewave}).

The pulse area
\begin{equation*}
    \theta_j \;\equiv\; \int_{-\infty}^{\infty}\!\mathrm{d}t\,\Omega_j(t-t_j)
\end{equation*}
sets the corresponding rotation angle.
In the idealised delta-pulse limit one replaces $\Omega_j(t-t_j)\to \theta_j\,\delta(t-t_j)$, so the pulse acts as an instantaneous unitary
\begin{equation*}
    \hat{U}_j
    \;=\;
\exp\!\left[-\frac{\mathrm{i}\,\theta_j}{2}\,\big(\cos\phi_j\,\hat{\tau}_x+\sin\phi_j\,\hat{\tau}_y\big)\right],
\end{equation*}
where $\hat{\tau}_{x,y,z}$ denote the Pauli matrices.

\begin{figure}[htbp]
    \centering
    \includegraphics[width=0.7\textwidth]{figures/app_bloch_phase_rotation.pdf}
    \caption{Bloch-sphere visualization of one ultrashort interaction with a two-level system. The initial state is $\ket{\mathrm{g}}$. Panels show examples for (a) $\theta=\pi/2$ and (b) $\theta=\pi$ with different carrier phase kicks $\phi$. The stars mark the final state after the pulse.}
    \label{fig:delta_phase_kick_bloch_rotation}
\end{figure}

Successive phase kicks therefore accumulate factors of the form $e^{\pm i\phi_j}$ on the optical coherences. This is the geometric counterpart of the phase factors used in the main text: phase cycling isolates a desired pathway class by Fourier projecting the measured signal with respect to $(\phi_1,\phi_2,\phi_3)$.
